% Общий стиль C4-диаграмм: гротеск, белый текст с тонкой обводкой,
% палитра C4-model. Подключается каждым fig-*.tex через \input.
\usepackage{fontspec}
\setmainfont{Roboto}
\usepackage{contour}
\contourlength{0.5pt}
\usepackage{tikz}
\usetikzlibrary{arrows, positioning, shadows, automata, arrows.meta, fit, calc}

% Палитра C4-model, затемнённая до уровня, на котором белый текст читается
\definecolor{c4person}{HTML}{08427B}
\definecolor{c4system}{HTML}{1168BD}
\definecolor{c4container}{HTML}{2E74B5}
\definecolor{c4component}{HTML}{4A90D9}
\definecolor{c4store}{HTML}{2A6187}
\definecolor{c4ext}{HTML}{7A7A7A}

% Заголовок узла: белый полужирный с тёмной обводкой
\newcommand{\ttl}[1]{\contour{black!80}{\textbf{#1}}}
% Тип узла в квадратных скобках
\newcommand{\typ}[1]{{\small [#1]}}

% Слой для рамок-контейнеров (виртуальных машин), чтобы заливка не перекрывала узлы
\pgfdeclarelayer{background}
\pgfsetlayers{background,main}
